% ============================================================================
% บทที่ 5 — แก้ไขเวลาสอบด้วยตนเอง
% ============================================================================
\mychapter{แก้ไขเวลาสอบด้วยตนเอง}

\noindent บทนี้ครอบคลุมการแก้ไขเวลาสอบด้วยตนเอง ตั้งแต่การดูตัวอย่างผลก่อนบันทึก
การยอมรับและการปฏิเสธข้อเสนอ วิชาที่เชื่อมโยงหรือแก้ไขไม่ได้
วิชาที่ยังไม่ได้จัด และผลของการจัดตารางใหม่ต่อการแก้ไข

\section{หลักการพื้นฐาน}

การแก้ไขเวลาสอบด้วยตนเองในระบบมีหลักการสำคัญดังนี้:

\begin{itemize}
  \item \textbf{การดูตัวอย่างก่อนบันทึก} เมื่อป้อนวันที่และเวลาใหม่ ระบบจะแสดง
        ตัวอย่างผลลัพธ์ให้ดูก่อน โดยการดูตัวอย่างยังไม่ใช่การบันทึก
  \item \textbf{การบันทึกและล็อกเวลา} หลังจากดูตัวอย่างและยอมรับแล้ว ต้องกด
        \texttt{บันทึกและล็อกเวลา} เพื่อบันทึก เวลาที่บันทึกจะถูกล็อกไว้
  \item \textbf{วิชาที่เชื่อมโยง} วิชาที่เชื่อมโยงแบบสอบตรงกันจะย้ายเวลาพร้อมกัน
  \item \textbf{วิชาที่แก้ไขไม่ได้} วิชาที่มีเวลาจากไฟล์ต้นทางหรือถูกบล็อก
        จากข้อมูลต้นทางจะแก้ไขไม่ได้
\end{itemize}

\section{ขั้นตอนการแก้เวลา}

\begin{steps}
  \item ไปที่หน้า \textbf{รายวิชา} เลือกประเภทการสอบ (\texttt{กลางภาค} หรือ
        \texttt{ปลายภาค}) ค้นหาวิชาที่ต้องการ แล้วกดปุ่ม \texttt{ดูรายละเอียด}
  \item ในหน้าต่างรายละเอียด เลือก \textbf{วันที่} จากช่องเลือกวันที่
        (แสดงรูปแบบตามภาษาของเบราว์เซอร์ เช่น \texttt{20/08/2026})
        และป้อน \textbf{เริ่ม} และ \textbf{สิ้นสุด}
        เช่น เวลา \texttt{13:00} ถึง \texttt{16:00}
  \item กดปุ่ม \texttt{ตรวจสอบการเปลี่ยนแปลง} เพื่อดูตัวอย่างผล
  \item หากผลถูกต้อง กดปุ่ม \texttt{บันทึกและล็อกเวลา}
\end{steps}

\begin{expected}
เมื่อข้อเสนอผ่านการตรวจสอบ ระบบแสดงข้อความ
\textbf{ไม่พบข้อขัดแย้งที่ขัดขวางการแก้ไข จะล็อกเวลา 1 วิชา}
และปุ่ม \texttt{บันทึกและล็อกเวลา} พร้อมใช้งาน เมื่อกดบันทึกแล้ว
ระบบแสดงข้อความ \textbf{บันทึกและล็อกเวลาแล้ว 1 รายการ}
ตารางรายวิชาแสดงเวลาใหม่พร้อมสถานะ \texttt{ล็อกเวลา}
เช่นเดียวกับปฏิทินและตารางกลุ่มนักศึกษา
\end{expected}

\manualfigure{docs/usage-report/screenshots/crops/F13.png}{0.72\textwidth}{%
  ผลการตรวจสอบที่ยอมรับ แสดงข้อความยืนยันว่าจะล็อกเวลา}

\begin{note}
การกด \texttt{ตรวจสอบการเปลี่ยนแปลง} เป็นเพียงการดูตัวอย่าง
การปิดหน้าต่างก่อนกด \texttt{บันทึกและล็อกเวลา} จะไม่เปลี่ยนเวลาสอบ
และหากแก้ไขช่องใดช่องหนึ่งหลังจากดูตัวอย่างแล้ว ต้องกด
\texttt{ตรวจสอบการเปลี่ยนแปลง} อีกครั้ง เพราะตัวอย่างเดิมไม่ครอบคลุมค่าที่แก้ไขล่าสุด
\end{note}

ช่วงเวลา \texttt{09:00--12:00} และ \texttt{13:00--16:00} ที่ตั้งไว้ในบทที่ 3
ใช้กับการจัดอัตโนมัติเท่านั้น การแก้ไขด้วยตนเองสามารถป้อนช่วงเวลาอื่นได้
เช่น \texttt{10:30--13:30} หากไม่ขัดแย้งกับวิชาอื่น

\needspace{9\baselineskip}
\section{กรณีที่ระบบไม่ยอมรับการแก้ไข}

เมื่อข้อเสนอขัดกับเวลาของวิชาอื่นหรือกฎของข้อมูลต้นทาง ระบบจะแสดงข้อความปฏิเสธ
ปุ่ม \texttt{บันทึกและล็อกเวลา} จะไม่ทำงาน และเวลาเดิมของวิชาไม่เปลี่ยนแปลง:

\begin{center}
\small
\begin{tabular}{L{4.2cm}L{4.6cm}L{5.2cm}}
\toprule
\textbf{สาเหตุ} & \textbf{อาการ} & \textbf{วิธีแก้} \\
\midrule
เวลาทับซ้อนกับวิชาอื่นในกลุ่มนักศึกษาเดียวกัน &
ข้อความปฏิเสธระบุวิชาที่ขัดแย้งและช่วงเวลาที่ทับซ้อน &
ปรับวันหรือเวลาใหม่แล้วกด \texttt{ตรวจสอบการเปลี่ยนแปลง} อีกครั้ง \\
วิชาสอบเต็มวัน &
ข้อเสนอครึ่งวันถูกปฏิเสธ &
ใช้ช่วงเต็มวันที่ตั้งค่าไว้ เช่น \texttt{09:00--16:00} \\
วิชามีเวลาจากไฟล์ต้นทางหรือถูกบล็อก &
ช่องวันที่และเวลาเป็นสีจางและแก้ไขไม่ได้ &
แก้ไขไฟล์ต้นทางแล้วนำเข้าใหม่ (ดูบทที่ 6) \\
\bottomrule
\end{tabular}
\end{center}

ตัวอย่างเช่น ป้อน \texttt{2026-08-21 09:00--12:00} สำหรับ \course{061000001}
ซึ่งมีวิชาอื่นของกลุ่ม \texttt{DEMO-A} สอบอยู่ช่วงนั้นแล้ว แล้วกด
\texttt{ตรวจสอบการเปลี่ยนแปลง} ระบบจะแสดงข้อความปฏิเสธ เช่น
\textbf{ไม่สามารถบันทึกการเปลี่ยนแปลงได้} ให้ปรับข้อเสนอแล้วดูตัวอย่างใหม่

\manualfigure{docs/usage-report/screenshots/crops/F15.png}{0.72\textwidth}{%
  ข้อเสนอที่ถูกปฏิเสธเนื่องจากเวลาสอบทับซ้อน}

\begin{caution}
บทบาทของไฟล์ (\texttt{จัดอัตโนมัติ} หรือ \texttt{จัดไปแล้ว})
ไม่ได้เป็นตัวกำหนดว่าแต่ละวิชาจะแก้ไขได้หรือไม่
วิชาที่มีเวลาจากไฟล์ต้นทางจะแก้ไขไม่ได้เสมอ
\end{caution}

วิชาที่มีเวลามาจากไฟล์ต้นทางหรือถูกบล็อกจากข้อมูลต้นทางจะแก้ไขไม่ได้
ช่องวันที่และเวลาในหน้าต่างรายละเอียดจะเป็นสีจางและกดไม่ได้ เช่น
\course{061000006} ซึ่งอยู่ในไฟล์ \texttt{จัดอัตโนมัติ} แต่มีเวลาจากไฟล์ต้นทาง

\manualfigure{docs/usage-report/screenshots/crops/F18.png}{0.72\textwidth}{%
  วิชาที่มีเวลาจากไฟล์ต้นทาง \course{061000006} ช่องวันที่และเวลาแก้ไขไม่ได้}

\section{การย้ายวิชาที่เชื่อมโยงแบบสอบตรงกัน}

วิชาที่กำหนดให้สอบตรงกันจะย้ายเวลาพร้อมกันทั้งกลุ่ม:

\begin{steps}
  \item เปิดรายละเอียดของ \course{061000003} (กลางภาค)
        ซึ่งเชื่อมโยงกับ \course{061000004}
  \item ป้อนวันที่ \texttt{2026-08-21} เวลา \texttt{13:00} ถึง \texttt{16:00}
  \item กด \texttt{ตรวจสอบการเปลี่ยนแปลง}
\end{steps}

\begin{expected}
ระบบแจ้งจำนวนวิชาที่จะย้ายพร้อมกัน เช่น \textbf{จะล็อกเวลา 2 วิชา}
เมื่อกด \texttt{บันทึกและล็อกเวลา} ทั้งสองวิชาย้ายไปเวลาเดียวกัน
และถูกบันทึกเป็น \texttt{ล็อกเวลา}
\end{expected}

\manualfigure{docs/usage-report/screenshots/crops/F16.png}{0.72\textwidth}{%
  การย้ายวิชาที่เชื่อมโยงแบบสอบตรงกัน แสดงจำนวนวิชาที่จะย้ายพร้อมกัน}

\section{การแก้ไขเวลาที่ล็อกไว้แล้วซ้ำ}

วิชาที่ล็อกเวลาไว้แล้วสามารถแก้ไขได้อีก ด้วยขั้นตอนเดียวกับปกติ
คือเปิดรายละเอียด ป้อนเวลาใหม่ ดูตัวอย่าง แล้วกด \texttt{บันทึกและล็อกเวลา}
เวลาของวิชาจะเปลี่ยนเป็นค่าใหม่ และสถานะยังคงเป็น \texttt{ล็อกเวลา}

\section{การกำหนดเวลาวิชาที่ยังไม่ได้จัด}

หากมีวิชาที่ยังไม่ได้จัด (เช่น จากปัญหาข้อมูลต้นทาง) ระบบจะแสดงจำนวน
ในหน้าภาพรวม ผู้ใช้กำหนดเวลาวิชาเหล่านี้เองได้:

\begin{steps}
  \item ไปที่หน้า \textbf{รายวิชา} และกรองสถานะเป็น \texttt{ยังไม่ได้จัดตาราง}
  \item เปิดรายละเอียดของวิชาที่ยังไม่ได้จัด
  \item ป้อนวันที่และเวลา กด \texttt{ตรวจสอบการเปลี่ยนแปลง}
        แล้วกด \texttt{บันทึกและล็อกเวลา}
\end{steps}

\begin{caution}
การกำหนดเวลาด้วยตนเองไม่สามารถเลี่ยงกฎของข้อมูลต้นทางได้
หากวิชาถูกบล็อกเนื่องจากข้อมูลต้นทางไม่สมบูรณ์ (เช่น กลุ่มนักศึกษาขาดหาย)
ต้องแก้ไขไฟล์ต้นทางแล้วนำเข้าใหม่ก่อน (ดูบทที่ 6)
\end{caution}

\section{การจัดตารางใหม่และผลต่อการแก้ไข}

สิ่งสำคัญที่ต้องเข้าใจคือความแตกต่างระหว่างการจัดตารางใหม่และการนำเข้าใหม่:

\begin{description}
  \item[\textbf{จัดตารางใหม่}] คำนวณตารางใหม่จากข้อมูลที่นำเข้าไว้เดิม
        การล็อกเวลาที่ผู้ใช้ทำไว้จะถูกยกเลิกทั้งหมด
  \item[\textbf{นำเข้าใหม่}] แทนที่โครงการที่นำเข้าไว้เดิมทั้งหมด ล้างผลลัพธ์
        ล้างการล็อกเวลา และตัดการเชื่อมต่อคลาวด์/รายการแชร์ปัจจุบัน
        การบันทึกครั้งแรกหลังนำเข้าใหม่จะสร้างโครงการคลาวด์ใหม่
\end{description}

\begin{caution}
การจัดตารางใหม่จะยกเลิกเวลาที่ล็อกด้วยตนเองจากผลลัพธ์ปัจจุบันเสมอ
ก่อนจัดตารางใหม่ ควรบันทึกโครงการขึ้นคลาวด์หากต้องการเก็บฉบับที่แก้ไขไว้ใช้งานต่อ
และสามารถดาวน์โหลด \file{project.zip} เก็บเป็นหลักฐานเพิ่มได้ (ดูบทที่ 7)
\end{caution}

\manualfigure{docs/usage-report/screenshots/figures/F19.png}{0.95\textwidth}{%
  สถานะก่อนการจัดตารางใหม่ วิชาที่ล็อกเวลาไว้แล้ว}

\manualfigure{docs/usage-report/screenshots/figures/F19b.png}{0.95\textwidth}{%
  สถานะหลังการจัดตารางใหม่ การล็อกเวลาถูกยกเลิก}

\seealso{บทที่ 6 สำหรับการแก้ปัญหาข้อมูลต้นทาง และบทที่ 7 สำหรับการเก็บงานและการเปิดงานต่อ}
